\documentclass[11pt]{article}
\usepackage[utf8]{inputenc}
\usepackage[T1]{fontenc}
\usepackage{lmodern}
\usepackage[spanish]{babel}
\usepackage{amsmath,amssymb}
\usepackage{graphicx}
\usepackage[left=2cm,right=2cm,top=2cm,bottom=2cm]{geometry}
\usepackage{enumitem}
\usepackage{booktabs}
\usepackage{array}
\usepackage{xcolor}
\usepackage{framed}

\spanishdecimal{.}

\setlist[enumerate,1]{label=\textbf{\arabic*.}, leftmargin=*, itemsep=6pt}
\setlist[enumerate,2]{label=(\alph*), itemsep=4pt}
\setlength{\parindent}{0pt}
\setlength{\parskip}{4pt}
\pagestyle{plain}

\newcommand{\Ein}{E_{\text{in}}}
\newcommand{\Eout}{E_{\text{out}}}
\newcommand{\Hyp}{\mathcal{H}}
\newcommand{\dvc}{d_{\text{VC}}}
\newcommand{\vx}{\mathbf{x}}
\newcommand{\vw}{\mathbf{w}}
\newcommand{\vy}{\mathbf{y}}
\newcommand{\R}{\mathbb{R}}
\newcommand{\T}{^{\top}}
\newcommand{\E}{\mathbb{E}}
\newcommand{\Prob}{\mathbb{P}}
\newcommand{\pts}[1]{\hfill\textbf{(#1~puntos)}}
\newcommand{\linea}[1]{\rule{#1}{0.4pt}}
\newcommand{\ptsinc}[1]{{\unskip\nobreak\hfil\penalty50\hskip1em\hbox{}\nobreak\hfill(#1)\parfillskip=0pt\finalhyphendemerits=0\par}}

\newcommand{\encabezado}[2]{%
\begin{center}
{\large\bfseries Reconocimiento de Patrones}\\[2pt]
Universidad de Sonora \textperiodcentered{} Licenciatura en Ciencias de la Computación\\
Semestre 2027-1 \textperiodcentered{} Prof.\ Julio Waissman Vilanova\\[10pt]
{\Large\bfseries #1}%
\ifdefined\tipodoc\\[4pt]{\large\itshape\tipodoc}\fi
\end{center}
\ifdefined\tipodoc\else
\vspace{6pt}
\textbf{Nombre:} \hrulefill\quad \textbf{Expediente:} \rule{3.5cm}{0.4pt}\par
\vspace{6pt}
\fi
#2
\par\medskip\hrule\medskip
}

\newcommand{\espacio}[1]{\par\vspace{#1}}
\newcommand{\nuevapagina}{\newpage}


\begin{document}

\encabezado{Examen rápido 7 \textperiodcentered{} Series de tiempo}{%
Tiempo: 20 minutos. Responde en esta hoja. Puedes usar calculadora; no se permiten apuntes. Deja indicadas las operaciones.}

\begin{enumerate}

%% ---------------------------------------------------------------
\item \textbf{Un AR(1).} La serie sigue $y_t = 2 + 0.6\, y_{t-1} + \varepsilon_t$, con $\varepsilon_t$ ruido blanco de varianza $\sigma^2 = 0.64$.
\begin{enumerate}
  \item ¿Es estacionaria? Calcula su media $\mu$, su varianza $\gamma_0$ y las autocorrelaciones $\rho_1$, $\rho_2$ y $\rho_3$.
  \espacio{3.2cm}

  \item El último dato es $y_T = 8$. Pronostica $y_{T+1}$ y $y_{T+2}$, y da un intervalo del 95\,\% para cada uno (supón errores gaussianos y parámetros conocidos).
  \espacio{3.5cm}

  \item ¿A qué tienden el pronóstico y el ancho del intervalo cuando el horizonte $h$ crece? ¿Qué cambiaría si $\phi = 1$?
  \espacio{2.2cm}
\end{enumerate}

\nuevapagina
%% ---------------------------------------------------------------
\item \textbf{Ingenuo, estacional y MASE.} Ventas trimestrales (en miles) de dos años de entrenamiento y del año de prueba:
\begin{center}
\begin{tabular}{l|cccc|cccc|cccc}
 & \multicolumn{4}{c|}{año 1} & \multicolumn{4}{c|}{año 2} & \multicolumn{4}{c}{prueba} \\
trimestre & 1 & 2 & 3 & 4 & 1 & 2 & 3 & 4 & 1 & 2 & 3 & 4 \\ \hline
ventas & 10 & 14 & 18 & 12 & 12 & 16 & 21 & 13 & 14 & 17 & 22 & 15
\end{tabular}
\end{center}
Con información hasta el final del año 2 se pronostican los cuatro trimestres de prueba.
\begin{enumerate}
  \item Escribe los pronósticos del método ingenuo y del estacional ingenuo ($m = 4$) y calcula el MAE de cada uno.
  \espacio{3cm}

  \item Calcula la escala $q = \frac{1}{T-m}\sum_{t > m} |y_t - y_{t-m}|$ con el entrenamiento, y el MASE de los dos métodos. ¿Qué significa un MASE menor que 1?
  \espacio{3cm}

  \item Un modelo nuevo obtiene MASE de 0.9 en estos cuatro trimestres de prueba. «Es menor que 1, así que le gana a la referencia estacional», dice su autor. ¿Tiene razón? ¿Qué modelo usarías?
  \espacio{2cm}
\end{enumerate}

%% ---------------------------------------------------------------
\item \textbf{Un AR(1) a mano, e identificación.}
\begin{enumerate}
  \item Con la serie $3, 5, 4, 6, 5$, ajusta $y_t = c + \phi\, y_{t-1}$ por mínimos cuadrados (usa los cuatro pares $(y_{t-1}, y_t)$) y pronostica el siguiente valor.
  \espacio{3.5cm}

  \item Para cada serie, propón un modelo: (i) la ACF decae como $0.8^k$ y la PACF solo tiene una barra, en el rezago 1; (ii) la ACF solo tiene barras en los rezagos 1 y 2 y la PACF decae; (iii) la ACF apenas baja: $0.99, 0.98, 0.97, \dots$
  \espacio{2.5cm}
\end{enumerate}

\end{enumerate}


\end{document}
